\section{Reverse induction and perfect-sequence extraction}\label{sec:induction}

At exposure step \(i\), append to the frozen one-step exponent
\[
-\mathsf d_{k,w,d}(P^*_{k,d})
-\widehat B_w(C)\frac{S_i}{d}
+\sum_{j>i}\frac{e_{ij}^2}{2\mu_wd}
+\frac{d^{-1/5}}{\sqrt d}\sum_{j>i}T_{ij}.
\tag{6.1}
\]
Every coefficient is fixed before the relevant column is sampled.  The
projection identity retains each future edge, and (4.6) supplies every
boundary edge with its full coefficient.

At the empty state, \cref{lem:multiplicity} yields
\[
P^*_{R,r}\le\mathcal H_C(r)
\exp\!\left[-\mathcal D_{C,w,d}(r)
-\frac1d\left\{2\widehat B_wA_0\binom r4+E_wK_r\right\}
+J_{r,d}\right],
\tag{6.2}
\]
where
\[
0\le J_{r,d}\le
\frac{m_0+\theta/D}{1-\rho}d^{-1/5}\binom r2=o_C(r^2).
\tag{6.3}
\]

For perfect-sequence extraction, relabel every retained ordered set of size
\(q\) and rebuild the deterministic recursion with target size \(q\).
All feasibility inequalities improve because \(q/\sqrt d\le1/D\).  If
\(u\) vertices are deleted, restoring the new deficit at size \(\ell\) costs
at most
\[
\frac1d\left[
2\widehat B_wA_0u\binom\ell3
+E_w\left(u\binom\ell2+2u\binom\ell3\right)
\right]
\le\left[\frac{\widehat B_wA_0+E_w}{3D^2}+o_C(1)\right]\ell u.
\tag{6.4}
\]
The bracket is \(O(1/\log C)<1\).  The unchanged source failure factor
\((p/10)^{10\ell u}\) dominates (6.4), the old extraction costs and (6.3).
Summing over deleted subsets remains \(o(1)\).  Thus the same source
first-moment argument applies with the new exponent intact.

